\scanpage{202}
\mspage{118}
% archive sheet number: 201
% element: text | top, lines 1-4 | grandes devant pi0
grandes devant $\pi_0$, cela implique que $\mathrm{Filt}^{\pi}$ est stable par facteurs directs, d'où $F\in\mathrm{Ob}\,\mathrm{Filt}^{\pi}$.

Enfin, si $M=\mathrm{Ens}$, $C=\unsure{\mathrm{Ens}(\leq\pi_0)|\mathrm{Ens}(\pi_0)}$. On note que
% element: display | upper third, line 5 | Filt pi equals set
\[
\mathrm{Filt}^{\pi}=\{X\in\mathrm{Ob}\,\mathrm{Ens}\mid\operatorname{card}(X)^{\pi_0}\leq\pi\}
\]
% element: text | upper middle, line 6 | D'autre part
D'autre part, on trouve facilement
% element: display | upper middle, lines 7-9 | Ens pi definition
\[
\begin{aligned}
(\mathrm{Ens})_\pi&\overset{\mathrm{def}}=\text{sous-catégorie pleine de }(\mathrm{Ens})\text{ formée}\\
&\qquad\text{des objets }\pi\text{-accessibles}\\
&=\{X\in\mathrm{Ob}\,\mathrm{Ens}\mid\operatorname{card}X\leq\pi\}
\end{aligned}
\]
% element: text | middle, lines 10-11 | Donc on a
Donc on a $(\mathrm{Ens})_\pi\subset\mathrm{Filt}^{\pi}$ ssi $\pi^{\pi_0}\leq\pi$, i.e. $\pi^{\pi_0}=\pi$, qed.
% element: margin note | right of conclusion, above corollary | Cf démonstration
\margin{Cf. \deleted{\unsure{faite}} \inserted{démonstration} de \unsure{la|(la} partie $1^\circ$) de 4.5.6, plus bas. (\unsure{fin de page} et p. suivante)}

% element: text | middle, line 12 | Corollaire 4.5.7
\uline{Corollaire 4.5.7}. Pour tt cardinal $\pi\geq\pi_1$, on a
% element: display | middle, line 13 | Filt pi subset M pi subset Filt
\[
\mathrm{Filt}^{\pi}\subset M_\pi\subset\mathrm{Filt}^{\pi^{\pi_0}}
\tag*{(4.5.7.1)}
\]
% element: text | lower half, lines 14-23 | L'inclusion
L'inclusion $\mathrm{Filt}^{\pi}\subset M_\pi$ est l'implication $a\Rightarrow c$ de la partie $1^\circ$ de 4.5.6. Pour prouver $M_\pi\subset\mathrm{Filt}^{\pi^{\pi_0}}$, \unsure{on|en} \deleted{étant noté que} \inserted{pose $\pi'=\pi^{\pi_0}$ et} \deleted{$\pi'$} on note que $\pi'{}^{\pi_0}=\pi^{\pi_0\cdot\pi_0}=\pi^{\pi_0}=\pi'$. Donc on peut appliquer $1^\circ$) à $\pi'$ \unsure{au|au lieu} $\pi$, et on trouve que $M_{\pi'}=\mathrm{Filt}^{\pi'}$. Comme $M_\pi\subset M_{\pi'}$, on en conclut $M_{\unsure{\pi|\pi'}}\subset\mathrm{Filt}^{\pi'}$, qed.

% element: text | foot, last three lines | J'ai oublié de démontrer
J'ai oublié de démontrer le $1^\circ$) dans 4.5.6. Vu le diagramme d'implications \unsure{en|(en} p.\,108, et \unsure{l'indication|l'implication} $a''\Rightarrow c$ déjà prouvée (p.\,109),
